\clearpage
\section{Representation Tracing: Coordinates, Controls, and Coverage}

\subsection{Fixed inputs and observer qualification}
Each reference set fixes input identities, order, context, masks, and coordinate correspondence across parameter states. Flattening an activation tensor retains every sample/token/spatial/feature entry. The cosine therefore compares changes on multiple \emph{matching} inputs, not averaged representations. The MLP reference set covers all 5,000 training images. ResNet uses the first intact original 1,000-image training-mode BatchNorm group. Transformer uses the first 64 prepared sequences and all 2,240 supervised positions. These latter reference sets are convenient fixed subsets, without replication across independently chosen reference sets or a claim of population-wide activation representativeness.

The observer preserves original model modes, RNG, buffers, normalization groups, and gradients. Identical-state repeat and plain-forward checks establish observer invariance in these deterministic settings. Training gradients still use the complete objective. The selected ResNet trace is the qualified native continuation around the 415 peak: all 187 loss states from 264 through 450 reproduce exactly, with initial gradient and state-296 parameter/gradient audit agreement. All six displayed site cosine/RMS histories agree with the paired motion figure. This qualifies that continuation; it is not bitwise recovery of a different runtime's original overview trajectory.

The measured endpoints are: MLP first linear map, first tanh, second linear map, second tanh, and output; ResNet first convolution, residual block groups 1/2/3, pooling, and output log probabilities; Transformer embedding, attention output and block output in each of two layers, and output log probabilities. These are computational boundaries, not averages over every internal operation. In particular, an attention-output trace does not separately identify query, key, value, and attention-weight mechanisms.

\subsection{Output spaces and loss relevance}
The MLP uses raw outputs, matching its half-squared error. ResNet and Transformer use all-class/all-vocabulary log probabilities, removing the softmax invariance to a common logit shift. Their Euclidean geometry still gives all coordinates equal weight; it is not automatically a loss-sensitivity metric. A correct-label log-probability readout has negative mean equal to the reference-set cross-entropy under the stated image/token weighting. This checks the loss readout without replacing the full output tensor in the primary SVD.

\subsection{Temporal reduction and descriptive controls}
Primary windows contain incoming changes at completed steps 451--501, 390--450, and 553--601, giving 51/61/49 rows. The extra predecessor used for the first cosine is excluded from the fitted matrix. The temporal Gram is accumulated in FP64 over coordinate blocks without example averaging, centering, random projection, or primary row normalization. Numerical PSD, trace, eigenvalue, and coefficient reconstruction checks qualify the reduction. Spatial patterns remain fixed over the fitted window; this does not imply that local dominant directions cannot rotate within their span.

For temporal Gram $K$, the orthogonal reference $K_{\rm orth}=\operatorname{diag}(\operatorname{diag}K)$ keeps all individual step norms and removes directional cross-terms. Unit-row normalization uses $K_{\rm unit}=B^{-1/2}KB^{-1/2}$ with $B=\operatorname{diag}(\operatorname{diag}K)$, excluding zero rows. Table~\ref{tab:t-spectral} separates amplitude weighting from directional structure. These references are not randomized significance tests. Since $PKP^\top$ has the same eigenvalues for every temporal permutation $P$, spectral concentration alone cannot establish oscillation; the signed coefficient series supplies temporal order.

\begin{table}[H]
\centering\small
\caption{Output-change spectral controls. $r_{95}$ is the number of modes needed for 95\% energy; $T$ is the fitted row count. Primary, orthogonal, and unit-row results use the same selected window.}
\label{tab:t-spectral}
\setlength{\tabcolsep}{6pt}
\begin{tabular}{@{}lr rr rr rr@{}}
\toprule
 & & \multicolumn{2}{c}{Observed} & \multicolumn{2}{c}{Same norms, orthogonal} & \multicolumn{2}{c}{Unit rows} \\
Case & $T$ & Top 1 (\%) & $r_{95}$ & Top 1 (\%) & $r_{95}$ & Top 1 (\%) & $r_{95}$ \\
\midrule
MLP & 51 & 53.02 & 5 & 7.85 & 42 & 36.21 & 5 \\
ResNet & 61 & 27.26 & 20 & 9.95 & 31 & 10.85 & 33 \\
Transformer & 49 & 41.65 & 15 & 16.72 & 20 & 16.77 & 29 \\
\bottomrule
\end{tabular}
\end{table}

Reference sets and windows differ across models; these are descriptive within-window measurements, rather than a controlled comparison of architecture dimensionality or causal propagation.

\subsection{Representation displacement retention}
For the same incoming changes $a{:}b$, let $K_\ell=D_\ell D_\ell^\top$.
Without truncating its modes, $R_{a:b}^\ell=\sqrt{\mathbf1^\top K_\ell\mathbf1}/\sum_i\sqrt{(K_\ell)_{ii}}$.
The direct Gram sum and full eigendecomposition agree numerically. Output values are
$.0170537/.0243649/.0350844$ for MLP/ResNet/Transformer; their first measured
sites give $.0325244/.0372071/.2433357$. The windows contain 51/61/49 changes.
A zero path length leaves retention undefined. Low retention alone need not
imply reversal: equal-length orthogonal changes yield $1/\sqrt{T_w}$.
